% https://easychair.org/conferences/?conf=ecai2025tutorials
% at most 4 pages

\documentclass{ecai} 

\usepackage{xspace}
\usepackage{natbib}
\usepackage{lstdoc}

\newcommand{\quotes}[1]{`#1'}
\newcommand{\iit}{Institute of Informatics and Telecommunications, NCSR~\quotes{Demokritos}\xspace}

 
\begin{document}

\begin{frontmatter}


\title{Introduction to Kolmogorov-Arnold Networks}

\author{\fnms{Tatiana}~\snm{Boura}}
\author{\fnms{Stasinos}~\snm{Konstantopoulos}}

\address{Institute of Informatics and Telecommunications,
  NCSR~\quotes{Demokritos}, Ag.~Paraskevi, Greece}


\begin{abstract}
The tutorial will introduce the attendees to Kolmogorov-Arnold
Networks (KANs), a new machine learning framework that promises to be
more interpretable and controllable than the multilayer perceptron (MLP).
Besides defining KANs and explaining how they are trained, the
tutorial will also guide attendees (through concrete examples) into
understanding how to interpret and manually control a trained KAN.  The
tutorial will also present the current state of the art regarding the
different architectures proposed within the framework and how they
have been applied to different tasks. As KANs are an innovative and
potentially impactful advancement in AI, we aim to help attendees
thoroughly understand their advantages and limitations and to motivate
the community to experiment with KANs in different domains and tasks. 
\end{abstract}

\end{frontmatter}

\hfill \break
\textbf{Proposed format:} Quarter day
%    Proposed format of the tutorial: half-day or quarter-day
\\
\textbf{Expected number of attendees:} 30--50



\section{Tutorial Topic}

%A brief description of why the tutorial topic would be of interest to
%a substantial part of the ECAI audience.

\paragraph{Objectives:}

\begin{itemize}
\item The tutorial aims to
  \emph{introduce expert non-specialists to an AI-Deep Learning subarea},
  that of Kolmogorov-Arnold Networks (KAN).
%
\item KANs are a brand new topic in machine learning, with the
  original publication of the idea published less than a year
  ago.\footnote{Z.~Liu, Y.~Wang, et al., KAN: Kolmogorov-Arnold
    Networks. arXiv:2404.19756v1 [cs.LG], Apr 2024.
    \url{https://arxiv.org/abs/2404.19756v1}}
  But it is a topic of \emph{emerging importance for AI}, since
  in this short time it has attracted impressive
  interest from the AI community with more than 10 extensions and
  applications published during 2024\footnote{Cf. Section~1 of
    D.~Basina, J.~R. Vishal, et al., 
    KAT to KANs: A review of Kolmogorov-Arnold Networks and the
    neural leap forward. arXiv:2411.10622 [cs.LG], Nov 2024.
    \url{https://arxiv.org/abs/2411.10622}}.
  %
\item The tutorial will \emph{explain} KANs, how they differ from the
  multi-layer perceptron, and what are the characteristics of machine
  learning tasks where they appear to be more promising,
  \emph{motivating} researchers to experiment with them.
  %
\item One of the core promises of the KAN framework is that its
  internal representation is more human-interpretable and more
  human-controllable than deep neural networks, so motivating the
  community to experiment with them can bring
  \emph{significant social impact} through advancing interpretable
  and controllable AI.
\end{itemize}  

\paragraph{Target Audience:}
% Who is the target audience of the tutorial (first-year PhD students?,
% experienced researchers?), including prerequisite knowledge?

The target audience is machine learning researchers and practitioners,
mostly those working on deep neural networks, regardless of seniority.
An understanding of the foundational concepts of machine learning and
artificial neural networks will be helpful, but a general
understanding of AI concepts should suffice.


\paragraph{Knowledge Conveyed:}
%What knowledge will be conveyed?
% What techniques/methods, concepts, or modeling frameworks will be conveyed
%Will it be demonstrated in particular application scenarios? If so, in which?

\begin{itemize}
\item The definition and core concepts behind Kolmogorov-Arnold
  Networks (KAN), and new machine learning framework.
\item The relationship between the KAN framework and the
  Kolmogorov-Arnold Theorem (KAT), as well as previous
  work on the relationship between KAT and machine learning.
\item The method for training KANs from data and an understanding of
  the convergence proof and what this proof means for practical
  applications.
\item The different KAN architectures 
  proposed in the literature and the empirical results from their
  application.
\item The interpretability and controllability claims of the
  framework and hands-on examples of how to read and modify a trained
  KAN.
\end{itemize}



\paragraph{Innovative and relevant nature of the topic}
%Why is the topic innovative/relevant?

Kolmogorov-Arnold Networks are a new approach to trainable networks
where the parameters are not weights, but spline control points that
define the functions that should be applied to the outputs of the
previous layer. Removing weights from the representation has a
profound effect on the \emph{interpretability} and
\emph{controllability} of the network, since the internal
representation can now be described in terms of mathematical functions
connected into a \emph{crisp} network of composition and addition
operations.
%
Naturally, this comes at the price of flexibility, so to fully exploit
its potential we aim to motivate the AI community to experiment with
KANs, map their scope of application, and find ways to bring out their
advantages while alleviating their limitations.


\section{Outline}

\begin{enumerate}
\item Introduction and background [15min]
  \begin{itemize}
  \item Introduction and tutorial outline
  \item Kolmogorov-Arnold Theorem (KAT):
    Formulation and proof sketch, at the level of abstraction
    needed to facilitate understanding of how KANs operate
  \item KAT used as a tool for understanding deep learning
    and related results prior to Kolmogorov-Arnold Networks
  \end{itemize}
\item Kolmogorov-Arnold Networks (KAN) [30min]
  \begin{itemize}
  \item Definition of KAN and relationship to KAT
  \item Training procedure and sketch of the convergence proof
  \item Interpretability and controllability: What can be read from a
    trained network and how can a trained network be manually
    modified. Specific examples of reading out of a network and manually
    modifying it.
  \end{itemize}
  \filbreak
\item Specific KAN architectures and applications [30min]
  \begin{itemize}
  \item Experimental results using the original KAN architecture 
  \item Recurrent architectures for sequence processing
  \item Architectures for solving differential equations and other
    physics tasks
  \end{itemize}
\item Closing Remarks [15min]
  \begin{itemize}
  \item Positioning in the machine learning space:
    relationship to MLPs and to symbolic regression
  \item Advantages and limitations, scope of potential applications
  \item Ambitious and interesting future work items selected from the
    literature
  \end{itemize}
\end{enumerate}
    

\section{Presenters}
\noindent
\textbf{Tatiana Boura}\footnote{Email: tatianabou@iit.demokritos.gr,
  ORCID: \href{https://orcid.org/0009-0008-0656-4372}{0009-0008-0656-4372}}
BSc in Informatics and Telecommunications (Athens, 2022),
MSc in Artificial Intelligence (Piraeus, 2024).
During her studies, she graduated with honors from her BSc and was awarded a Scholarship for Academic Excellence
for both academic semesters of the master's program. She was also awarded by the University of Piraeus
\footnote{Cf. \href{https://msc-ai.iit.demokritos.gr/el/announcements/vraveysi-apofoiton-toy-dpms-tehniti-noimosyni-apo-panepistimio-peiraios}{Graduate Awards 2024}}.
She has been affiliated with the \iit since September 2023.
During this time, she has worked on multiple cutting-edge AI-centered projects, focusing on the field of explainable
AI (xAI), with an emphasis on neurosymbolic artificial intelligence.
She served as a reviewer for ECAI 2024 and is a member of the Greek AI Society.

\noindent
\textbf{Stasinos Konstantopoulos}\footnote{Email: konstant@iit.demokritos.gr,
  ORCID: \href{https://orcid.org/0000-0002-2586-1726}{0000-0002-2586-1726}}
MEng in Computer Engineering and Informatics (Patras, 1997),
MSc in Artificial Intelligence (Edinburgh, 1998),
PhD on Inductive Logic Programming (Groningen, 2003), has been affiliated
with the \iit since 2005. During this time he has worked on various areas
of AI from computational logics to robot sensing to machine learning.
He has been a member of the Greek AI Society since 2006 where he was
PC Chair for its national AI conference (SETN 2010) and also served
on the society's 2021-2022 Board. He is a regular reviewer for AI and
Semantic Web conferences and journals, including ECAI.
He teaches \textit{Fundamentals and Background on AI},
\textit{Knowledge Representation and Reasoning} and
\textit{AI Applications} at the MSc on AI jointly organized by
University of Piraeus and NCSR \quotes{Demokritos}.\footnote{Cf.
  \url{https://msc-ai.iit.demokritos.gr/en/people}}

% Background in the tutorial area, including a list of publications/presentations
% example of work in the area (ideally, a published tutorial-level
% article or presentation materials on the subject)

\noindent
\textbf{Background:}
%
Tatiana and Stasinos' research on KANs focuses on recurrency and
sequence processing where they developed \emph{seqKAN}, a new KAN
architecture.\footnote{Submitted, under review.
  Pre-print: \url{https://arxiv.org/abs/2502.14681},
  implementation: \url{https://github.com/data-eng/seqKAN}
  and \url{https://pypi.org/project/seqkan}}
%
In November 2024 they co-organized and jointly presented a 2-hour
internal talk and workshop at the \iit with the aim to present the new
framework to colleagues, present their own work on the subject, and
motivate colleagues to apply KANs to different tasks.\footnote{Slides
and bibliography handout are available at
  \url{https://github.com/data-eng/KAN/tree/master/background}}
The event attracted around 40 participants.\footnote{Cf. 
  \url{https://www.iit.demokritos.gr/newsevents/iit-talk-on-kolmogorov-arnold-networks-theoretical-background-current-approaches-and-potential-next-steps}}


\newpage
\section{Information for communication purposes}

% A two-sentence description of the tutorial, suitable for inclusion
% in the conference brochure

\paragraph{Two-sentence description:}

The tutorial introduces attendees to Kolmogorov-Arnold Networks
(KANs), a new approach to machine learning that was recently proposed
as a more interpretable and controllable framework than the ubiquitous
multi-layer perceptron. During the tutorial attendees will learn what
the KAN framework is, how it is positioned in the general machine
learning landscape, what is meant by the framework's strong
interpretability and controllability claims, and what the literature has
shown about their efficacy and scope of application.


%A two-paragraph description of the tutorial, suitable for inclusion on
%the conference website. It should mention the target audience

\paragraph{Two-paragraph description:}

The tutorial introduces attendees to Kolmogorov-Arnold Networks
(KANs), a new approach to machine learning that was recently proposed
as a more interpretable and controllable framework than the ubiquitous
multi-layer perceptron. Besides providing an understanding of the KAN
framework and its position within the overall machine-learning
picture, the tutorial will also guide attendees through the different
KAN architectures that have been proposed in the literature and their
applications.

The target audience is machine learning researchers and practitioners,
and especially those working on deep neural networks. Besides the more
theoretical understanding of KAN architectures, the tutorial will also
present examples of how the interpretability and controllability of
KANs is concretely realized. Finally the tutorial will discuss
advantages and limitations of KANs for different application areas
with the aim to motivate interested members of the audience to
experiment with them.


%A brief bio statement of the presenter(s)
%For each presenter, a link to a picture that may be used by ECAI

\paragraph{Brief bio statement, Tatiana Boura:}

Tatiana (BSc in Informatics and Telecommunications, MSc in AI) is
affiliated to the \iit where she has worked on
cutting-edge AI-centered projects, focusing on the field of explainable
AI (xAI), with an emphasis on neurosymbolic artificial intelligence.
She is a member of the Greek AI Society.

\noindent
Photo: \url{https://avatars.githubusercontent.com/u/61378350}

\paragraph{Brief bio statement, Stasinos Konstantopoulos:}

Stasinos (MEng on Computer Engineering and Informatics, MSc on AI,
PhD on ILP) is affiliated to the \iit where he has worked on various areas
of AI from computational logics to robot sensing to machine
learning. He is a member of the Greek AI Society and served on the
society's 2021-2022 Board. He is a regular reviewer for AI and
Semantic Web conferences and journals, including ECAI 2025. He is also
an instructor at the MSc on AI jointly organized by University of
Piraeus and NCSR~\quotes{Demokritos}.

\noindent
Photo: \url{https://avatars.githubusercontent.com/u/663585}

\end{document}

% LocalWords:  Kolmogorov Tatiana Boura Stasinos Konstantopoulos KAT
% LocalWords:  perceptron ORCID BSc neurosymbolic ECAI MEng Patras
% LocalWords:  Groningen SETN seqKAN MSc interpretability controllability
% LocalWords:  xAI
